\documentclass[11pt]{article}
\usepackage[T1]{fontenc}
\usepackage[utf8]{inputenc}
\usepackage[margin=0.85in]{geometry}
\usepackage{amsmath,amssymb}
\usepackage{booktabs}
\usepackage{array}
\usepackage{longtable}
\usepackage{enumitem}
\usepackage{url}
\usepackage{xcolor}
\usepackage[colorlinks=true,linkcolor=blue,citecolor=blue,urlcolor=blue]{hyperref}
\setlength{\parindent}{0pt}
\setlength{\parskip}{0.48em}
\setlist[itemize]{leftmargin=1.4em,itemsep=0.16em,topsep=0.18em}
\setlist[enumerate]{leftmargin=1.6em,itemsep=0.16em,topsep=0.18em}
\newcommand{\clip}{\operatorname{clip}}
\newcommand{\sigm}{\sigma}
\title{ODD Protocol Description for the Agent-Based Model\\\large Generative AI as an Interaction Intermediary: Collaboration, Trust, and Social Learning}
\author{Yakup Turgut\\K\i rklareli University, K\i rklareli, Turkey}
\date{Repository model documentation}

\begin{document}
\maketitle

\begin{abstract}
This document provides a detailed Overview--Design concepts--Details (ODD) description of the agent-based model used in the study \emph{Generative AI as an Interaction Intermediary: Modeling How AI-Mediated Interaction Changes Collaboration, Trust, and Social Learning}. Heterogeneous agents repeatedly solve multidimensional tasks through peer-first, AI-first, or hybrid support. Their choices depend on learned task-specific expectations, source trust, local social information, and access or coordination costs. Realized outcomes update knowledge, confidence, source trust, route expectations, and weighted social ties. The document specifies the model entities, state variables, scheduling, behavioral equations, stochastic processes, reference parameter values, experimental design, output measures, and reproducibility procedures following the ODD protocol \cite{grimm2006standard,grimm2010odd,grimm2020odd} with decision-process clarification inspired by ODD+D \cite{muller2013odd}.
\end{abstract}

\tableofcontents
\newpage

\section{Purpose of this model document}
This document is the repository-level specification of the simulation model. Its purpose is to provide enough detail for an independent reader to understand, inspect, reproduce, and modify the model beyond the compressed description possible in the proceedings paper. It has four functions:
\begin{enumerate}
\item describe the model using the ODD structure;
\item make state variables, behavioral rules, parameter values, and scheduling explicit;
\item explain how abstract constructs such as knowledge, trust, AI access, peer expertise, and tie strength are operationalized;
\item connect the conceptual description to the Python implementation, experiment scripts, saved result files, and figures.
\end{enumerate}
The implementation and all repository documentation are available at \url{https://github.com/yakupturgut/generative-ai-interaction-abm}.

\section{Model scope and interpretation}
The model is an exploratory, mechanism-based ABM rather than a calibrated predictor of a particular workplace, classroom, platform, or population. It formalizes transparent assumptions about repeated problem solving when both generative AI and human peers can be used as support resources. System-level patterns are generated from local interactions and adaptive feedback, consistent with the generative logic of social simulation \cite{epstein2006generative}.

The constructs are operational. \emph{Knowledge} denotes task-relevant capability in a three-dimensional simulated requirement space. \emph{AI literacy} denotes the ability to extract productive value from AI output. \emph{Sociability} affects the accessibility and transfer efficiency of peer support. \emph{Verification} affects the ability to scrutinize AI output and integrate multiple sources. \emph{Trust} is a learned source-specific belief that moves toward experienced source quality. \emph{Tie strength} is a weighted representation of an active social relationship whose value changes through use and non-use.

The route-choice scores are bounded-rational behavioral propensities. They are not empirically estimated welfare utilities. Probabilistic softmax choice provides stochastic adaptation under heterogeneous expectations and imperfect information, drawing on random-utility ideas \cite{mcfadden1974conditional} and bounded rationality \cite{simon1955behavioral}.

\section{Overview}
\subsection{Purpose and patterns}
The model examines how repeated problem-solving choices among \textbf{Peer-first}, \textbf{AI-first}, and \textbf{Hybrid} support shape:
\begin{itemize}
\item the composition of AI use and human interaction;
\item task quality and task success;
\item domain-specific knowledge accumulation;
\item trust in AI and partner-specific peer trust;
\item social-tie reinforcement, decay, retention, and formation;
\item performance and knowledge inequality;
\item recovery after temporary or persistent degradation in AI reliability.
\end{itemize}
The central mechanism is a repeated feedback loop. A task arrives, the agent forms expectations for feasible support routes, a route is selected probabilistically, the task outcome is realized, and the resulting experience updates future expectations, trust, knowledge, confidence, and social relations.

\subsection{Entities, state variables, and scales}
The model contains human agents, tasks, a weighted undirected social network, and an AI support environment. There is no physical geography. One episode is an abstract collaborative problem-solving opportunity rather than a fixed unit of calendar time.

\begin{longtable}{p{0.16\textwidth}p{0.49\textwidth}p{0.27\textwidth}}
\caption{Principal agent-level state variables.}\label{tab:agentstates}\\
\toprule
Variable & Interpretation & Scale / role \\
\midrule
\endfirsthead
\toprule
Variable & Interpretation & Scale / role \\
\midrule
\endhead
$\mathbf{k}_i=(k_{i1},k_{i2},k_{i3})$ & Three-dimensional task-relevant skill vector of agent $i$. & $[0,1]^3$; changes through practice and novel information. \\
$a_i$ & AI literacy; ability to productively use AI output. & $[0,1]$; affects positive AI uptake and hybrid integration. \\
$s_i$ & Sociability; ability/orientation to engage effectively with peers. & $[0,1]$; affects peer cost, search, transfer, and learning. \\
$v_i$ & Verification tendency. & $[0,1]$; affects AI checking, hybrid integration, and harmful-output filtering. \\
$c_i$ & Confidence. & $[0,1]$; contributes to realized task quality and adapts slowly. \\
$A_i$ & Individual AI access/ease. & $[0,1]$; shapes AI availability and access friction. \\
$t_i^{AI}$ & Trust in AI. & $[0,1]$; moves toward experienced task-relevant AI source quality. \\
$t_{ij}^{P}$ & Agent $i$'s trust in peer $j$. & $[0,1]$; partner-specific and experience-based. \\
$V_{iam}$ & Expected realized quality of mode $m$ for task anchor $a$. & $[0.05,0.95]$; similarity-weighted experience memory. \\
$E_{ija}$ & Agent $i$'s expectation of peer $j$'s expertise for task anchor $a$. & $[0.01,0.99]$; partner- and task-specific. \\
$L_{iam}$ & Local observational signal for mode $m$ and task anchor $a$. & $[-0.18,0.24]$; updated from observed neighbors' marginal support gains. \\
$w_{ij}$ & Social-tie strength. & $[0,1]$; reinforced by useful peer contribution, decays when unused. \\
\bottomrule
\end{longtable}

Each task contains a three-dimensional requirement vector $\mathbf{r}$, ambiguity $u$, difficulty $d$, success threshold $h$, and a task-family label used for interpretation. The reference model contains $N=180$ agents and $T=150$ episodes.

\subsection{Process overview and scheduling}
At each episode, the following sequence is executed:
\begin{enumerate}
\item apply any scheduled change in AI reliability;
\item freeze the current skills and ties for within-episode reference;
\item generate one task for each agent;
\item draw AI availability for each agent and identify reachable peers;
\item compute route costs, expected values, trust terms, and local social signals;
\item convert feasible route propensities to probabilities and sample Peer-first, AI-first, or Hybrid;
\item if a peer is used, choose the partner probabilistically from perceived expertise, trust, tie strength, and search cost;
\item realize task quality through multidimensional task--resource matching;
\item classify task success using the task-specific threshold;
\item update route expectations, AI trust, peer trust, partner expertise expectations, knowledge, and confidence;
\item update used social ties from the peer's marginal contribution;
\item decay unused ties and remove ties below the removal threshold;
\item update local observational-learning memories from visible neighbors;
\item record episode-level and, when requested, task- and agent-level outputs.
\end{enumerate}
Agents are processed in a random permutation within each episode. Task generation is simultaneous, while state updates occur as agents complete their decisions.

\section{Design concepts}
\subsection{Basic principles}
The model combines five principles: heterogeneous capabilities, bounded probabilistic choice, multidimensional task--resource matching, experience-based learning, and co-evolving social relationships. The same route-choice architecture is used for Peer-first, AI-first, and Hybrid support. Differences in realized outcomes arise from source availability, task requirements, source capability, agent attributes, partner characteristics, and accumulated experience.

\subsection{Emergence}
Aggregate patterns such as AI-use rate, human-interaction rate, trust trajectories, knowledge inequality, and social-tie activity are system-level outcomes of repeated micro decisions. They are not direct state variables chosen by agents. In particular, AI-first use changes social structure through opportunity displacement: episodes spent without a peer provide no peer interaction with which to reinforce a social tie, while unused ties continue to decay slowly.

\subsection{Adaptation}
Agents adapt after realized outcomes. Route expectations are updated toward the experienced quality of the selected route for similar tasks. AI trust is updated toward experienced AI source quality. Peer trust and peer-expertise expectations are updated toward experienced partner quality. Confidence adapts toward task quality. Knowledge grows through practice and novel information. Social ties adapt through the realized marginal contribution of peer interaction.

\subsection{Objectives and decision-making}
Agents seek support for the current task rather than optimizing a long-run social objective. The route propensity for agent $i$ and route $m\in\{P,A,H\}$ is
\begin{equation}
U_{im}=\beta_Q\widehat Q_{im}+\beta_T T_{im}+\beta_L L_{im}-C_{im},
\label{eq:oddutility}
\end{equation}
where $\widehat Q_{im}$ is learned expected realized quality, $T_{im}$ is source trust, $L_{im}$ is local observational information, and $C_{im}$ is route cost. Feasible-route probabilities are
\begin{equation}
\Pr_i(m)=\frac{\exp(U_{im}/\tau)}{\sum_{r\in\mathcal F_i}\exp(U_{ir}/\tau)},
\end{equation}
with an exploration floor applied over feasible routes.

\subsection{Learning}
Three learning channels operate simultaneously:
\begin{enumerate}
\item \textbf{Experiential route learning}: selected-route quality updates task-similarity-weighted expected quality.
\item \textbf{Source learning}: source quality updates AI or partner-specific peer trust, and peer interaction updates perceived partner expertise.
\item \textbf{Capability learning}: task practice and novel support information increase the agent's three-dimensional skill vector.
\end{enumerate}

\subsection{Prediction}
Agents use task-anchor memories to predict route quality for the current task. Current task requirements are mapped to the three anchors through smooth cosine-similarity weights. Peer expertise is perceived imperfectly. The model therefore separates latent source capability from the agent's current expectation of that capability.

\subsection{Sensing}
Each agent has access to its own traits, current task requirements, personal route memories, AI trust and access, available peers, partner-specific peer expectations and trust, and current tie strengths. A stochastic subset of neighbors' recent support outcomes is observed through the local social-learning mechanism. Agents observe only the outcome of the route they actually use.

\subsection{Interaction}
Peer-first and Hybrid routes involve another human agent. Peer availability depends on current network neighbors, peer-response probability, and occasional external friend-of-friend search. A selected peer transmits task-relevant capability imperfectly. Social interaction can reinforce an existing tie or create a new tie when a newly discovered peer provides sufficiently positive marginal value.

\subsection{Stochasticity}
Stochasticity enters through:
\begin{itemize}
\item heterogeneous initial traits and skill dimensions;
\item initial social-network topology and tie weights;
\item task family, task requirement vector, ambiguity, and difficulty;
\item AI capability noise;
\item peer communication noise;
\item AI availability and peer response;
\item external peer discovery;
\item probabilistic route choice and partner choice;
\item observation of neighbors;
\item final task-quality noise.
\end{itemize}
Independent replications use deterministic seed construction so that paired experimental contrasts can use common random numbers.

\subsection{Collectives}
The social network constitutes a dynamic collective structure. Network-level quantities include density, weighted degree, active-tie fraction, tie retention, new-edge ratio, strong-tie prevalence, and unique peer partners. Collective interaction composition is measured through Peer-first, AI-first, and Hybrid shares.

\subsection{Observation}
Two overlapping interaction indicators are used:
\begin{align}
\text{Human interaction rate} &= \text{Peer-first share}+\text{Hybrid share},\\
\text{AI-use rate} &= \text{AI-first share}+\text{Hybrid share}.
\end{align}
Because Hybrid includes both AI and another person, these indicators are not complements and need not sum to one.

\section{Details}
\subsection{Initialization}
The initial skill profile combines a general capability component with domain-specific deviations. General capability is drawn from a bounded normal distribution centered near the reference knowledge mean; three independent domain deviations generate $\mathbf{k}_i$. AI literacy, sociability, verification, confidence, and AI access are bounded normal traits. Initial task-specific route values are centered near 0.50 with small random variation. AI trust is centered near 0.50.

The initial social graph is a connected Watts--Strogatz small-world network \cite{watts1998collective}. The nearest-neighbor parameter is selected to approximate the target density. Initial tie weights are uniformly drawn from $[0.42,0.68]$. Partner-specific peer trust is centered near 0.50. Initial peer-expertise expectations combine a neutral prior with noisy information about the connected partner's skill profile.

\subsection{Reference parameterization}
Tables~\ref{tab:param_ai}--\ref{tab:param_update} report the full reference parameterization stored in \texttt{ModelConfig} and \texttt{results/analysis\_configuration.json}.

\begin{longtable}{p{0.37\textwidth}p{0.15\textwidth}p{0.38\textwidth}}
\caption{Scale and AI-environment parameters.}\label{tab:param_ai}\\
\toprule Parameter & Value & Interpretation \\
\midrule\endfirsthead\toprule Parameter & Value & Interpretation \\
\midrule\endhead
Number of agents & 180 & Population size. \\
Episodes & 150 & Repeated task opportunities per run. \\
AI enabled & True & AI support available in the reference system. \\
AI reliability & 0.78 & Scales the realized AI capability profile. \\
Mean AI access & 0.60 & Population mean of individual AI ease/access. \\
AI access SD & 0.16 & Heterogeneity in individual AI access. \\
Verification support & 0.55 & Environmental support for checking/integrating AI output. \\
AI availability floor & 0.70 & Minimum availability under the saturating access mapping. \\
AI availability ceiling & 0.97 & Maximum availability under the saturating access mapping. \\
AI access saturation & 2.0 & Curvature of the access-to-availability mapping. \\
AI capability profile & $(0.94,0.64,0.76)$ & Baseline capability across three task dimensions. \\
AI capability noise SD & 0.085 & Episode-level capability noise. \\
AI output concentration & 22.0 & Configuration value retained for AI-output dispersion. \\
AI ambiguity penalty & 0.16 & Reduction in effective AI capability with task ambiguity. \\
\bottomrule
\end{longtable}

\begin{longtable}{p{0.37\textwidth}p{0.15\textwidth}p{0.38\textwidth}}
\caption{Peer, network, task, and agent-trait parameters.}\label{tab:param_peer}\\
\toprule Parameter & Value & Interpretation \\
\midrule\endfirsthead\toprule Parameter & Value & Interpretation \\
\midrule\endhead
Network density & 0.055 & Target initial small-world density. \\
Rewiring probability & 0.10 & Watts--Strogatz rewiring probability. \\
Peer-response probability & 0.80 & Probability an existing neighbor is reachable. \\
External peer-search probability & 0.12 & Base probability of outside-network search. \\
Peer-expertise visibility & 0.30 & Weight of noisy expertise information in the initial partner belief. \\
Social observation rate & 0.45 & Probability of observing a neighbor's recent outcome. \\
Peer communication noise SD & 0.095 & Noise in communicated peer capability. \\
Peer ambiguity penalty & 0.10 & Ambiguity-related reduction in peer transfer efficiency. \\
Minimum peer transfer efficiency & 0.28 & Lower bound on transferred peer capability. \\
Maximum peer transfer efficiency & 0.80 & Upper bound on transferred peer capability. \\
Task mix & $(1/3,1/3,1/3)$ & Structured, contextual, integrative family probabilities. \\
Mean task difficulty & 0.58 & Mean of the Beta difficulty draw. \\
Difficulty concentration & 10.0 & Concentration of the Beta difficulty distribution. \\
Task-similarity temperature & 0.090 & Softness of task-to-anchor similarity weights. \\
Mean initial knowledge & 0.52 & Center of general initial capability. \\
Mean AI literacy & 0.50 & Population AI-literacy center. \\
Mean sociability & 0.55 & Population sociability center. \\
Mean verification & 0.52 & Population verification center. \\
Mean confidence & 0.50 & Population confidence center. \\
Trait SD & 0.14 & Standard deviation for bounded heterogeneous traits. \\
\bottomrule
\end{longtable}

\begin{longtable}{p{0.37\textwidth}p{0.15\textwidth}p{0.38\textwidth}}
\caption{Choice and route-cost parameters.}\label{tab:param_choice}\\
\toprule Parameter & Value & Interpretation \\
\midrule\endfirsthead\toprule Parameter & Value & Interpretation \\
\midrule\endhead
Expected-quality weight $\beta_Q$ & 3.60 & Weight on learned route-quality expectation. \\
Trust weight $\beta_T$ & 0.12 & Weight on source trust. \\
Local social-learning weight $\beta_L$ & 0.55 & Weight on observed local support value. \\
Choice temperature $\tau$ & 0.31 & Softmax stochasticity. \\
Exploration floor & 0.025 & Minimum exploration mass per feasible route. \\
Peer base cost & 0.22 & Base peer-coordination cost. \\
AI base cost & 0.20 & Base AI-access cost before access saturation. \\
AI ambiguity verification cost & 0.14 & Extra checking burden on ambiguous AI use. \\
Hybrid extra cost & 0.12 & Additional integration/coordination cost. \\
Weak-network search cost & 0.10 & Search cost for low-strength or external peer access. \\
\bottomrule
\end{longtable}

\begin{longtable}{p{0.37\textwidth}p{0.15\textwidth}p{0.38\textwidth}}
\caption{Learning, network-update, and outcome parameters.}\label{tab:param_update}\\
\toprule Parameter & Value & Interpretation \\
\midrule\endfirsthead\toprule Parameter & Value & Interpretation \\
\midrule\endhead
Mode-value learning rate & 0.32 & Adaptation of task-specific route-quality expectations. \\
Trust learning rate & 0.12 & Adaptation of AI and peer trust. \\
Local-memory learning rate & 0.13 & Adaptation of observed neighborhood route signals. \\
Confidence learning rate & 0.030 & Adaptation of confidence toward realized quality. \\
Practice learning rate & 0.008 & Baseline learning from task practice. \\
Knowledge learning rate & 0.095 & Additional learning from novel support information. \\
Tie reinforcement rate & 0.085 & Positive update from useful peer contribution. \\
Negative tie-update rate & 0.040 & Negative update from harmful peer contribution. \\
Unused-tie decay & 0.0020 & Per-episode multiplicative decay of unused active ties. \\
Tie-removal threshold & 0.050 & Tie deleted below this strength. \\
Initial tie range & $[0.42,0.68]$ & Uniform initial active-edge weights. \\
New-tie initial strength & 0.34 & Starting strength after successful external contact. \\
New-tie formation threshold & 0.010 & Minimum positive peer marginal contribution for new tie. \\
Support uptake scale & 0.72 & Scale for beneficial source uptake. \\
Harmful-advice uptake scale & 0.40 & Scale for uptake of negative source information. \\
Hybrid coordination burden & 0.055 & Burden from combining AI and peer support. \\
Hybrid conflict scale & 0.20 & Penalty for disagreement between AI and peer information. \\
Quality noise SD & 0.024 & Stochastic noise added to final task quality. \\
Reference seed & 20260822 & Base configuration seed recorded with results. \\
\bottomrule
\end{longtable}

\subsection{Task generation}
The three task prototypes are
\begin{align}
\mathbf{p}_{S}&=(0.66,0.22,0.12),\\
\mathbf{p}_{C}&=(0.22,0.63,0.15),\\
\mathbf{p}_{I}&=(0.34,0.29,0.37),
\end{align}
with ambiguity centers $(0.24,0.66,0.48)$. Each agent receives a task family according to the task mix. Conditional on family $f$, the requirement vector is drawn from a Dirichlet distribution with concentration vector
\begin{equation}
\boldsymbol\alpha_f=7\mathbf{p}_f+0.70,
\end{equation}
then normalized by construction. Ambiguity is a bounded normal draw around the corresponding family center with SD 0.11. Difficulty is drawn from a Beta distribution with mean 0.58 and concentration 10, bounded to $[0.08,0.96]$. The task-specific success threshold is
\begin{equation}
h=\clip(0.46+0.22d+0.06u,\,0.48,\,0.76).
\end{equation}

Task similarity to the three anchors is based on cosine similarity and a softmax with temperature 0.090. These continuous weights determine how strongly experience on the current task updates each anchor memory.

\subsection{AI availability and access cost}
For individual access $A_i\in[0,1]$, the saturating access transformation is
\begin{equation}
s(A_i)=\frac{1-\exp(-kA_i)}{1-\exp(-k)},\qquad k=2.0,
\end{equation}
and AI availability probability is
\begin{equation}
p_i^{AI}=p_{\min}+(p_{\max}-p_{\min})s(A_i),
\end{equation}
with $p_{\min}=0.70$ and $p_{\max}=0.97$. AI usage friction is
\begin{equation}
C_i^{AI}=c_A(1-A_i)^{1.65}+c_u u(1-v_i),
\end{equation}
where $c_A=0.20$ and $c_u=0.14$.

\subsection{Peer availability, search, and partner selection}
Each existing neighbor independently responds with probability 0.80. External search occurs with probability
\begin{equation}
p_i^{search}=0.12(0.45+0.55s_i).
\end{equation}
External candidates are weighted by friend-of-friend connectivity when such information is available.

For a candidate peer $j$, partner-selection propensity is
\begin{equation}
Z_{ij}=0.50\widehat E_{ij}+0.25t_{ij}^{P}+0.25w_{ij}-C_{ij}^{search},
\end{equation}
where $\widehat E_{ij}$ is task-weighted perceived expertise. A softmax with temperature $\max(0.20,0.78\tau)$ determines the sampled partner.

\subsection{Route choice}
Peer route cost is
\begin{equation}
C_i^{P}=c_P(1-s_i)+c_W(1-w_i^{\max}),
\end{equation}
where $c_P=0.22$, $c_W=0.10$, and $w_i^{\max}$ is the strongest currently reachable tie. Hybrid integration ability for route-cost purposes is
\begin{equation}
I_i^{C}=\clip(0.36v_i+0.24a_i+0.22s_i+0.18V,0,1),
\end{equation}
with environmental verification support $V=0.55$. Hybrid cost is
\begin{equation}
C_i^{H}=0.52C_i^{P}+0.52C_i^{AI}+0.12(1-I_i^{C}).
\end{equation}
These costs enter Eq.~\ref{eq:oddutility}. Trust for Hybrid is the mean of AI trust and the prospective peer trust. Route expectations and local social signals are blended across task anchors using current task-similarity weights.

\subsection{AI support}
The reference AI capability profile is $\mathbf{g}=(0.94,0.64,0.76)$. At current reliability $R$, the expected capability vector moves between an uninformative floor 0.12 and the profile:
\begin{equation}
\bar{\mathbf g}=0.12+R(\mathbf g-0.12).
\end{equation}
Ambiguity scales the vector by $(1-0.16u)$, after which Gaussian noise with SD 0.085 is added and values are clipped to $[0.01,0.99]$.

AI information affects the focal agent's temporary effective skill through beneficial and harmful uptake. Beneficial uptake is scaled by AI literacy and verification; harmful uptake is filtered by verification and environmental verification support. Thus source quality and source use are represented separately.

\subsection{Peer support}
For peer $j$, communication efficiency is based on focal sociability, peer sociability, current tie strength, and partner-specific trust:
\begin{equation}
r_{ij}=0.29s_i+0.29s_j+0.24w_{ij}+0.18t_{ij}^{P}.
\end{equation}
This score is mapped to the interval $[0.28,0.80]$, reduced by task ambiguity, and used to move communicated peer capability from the focal agent's own capability toward the peer's latent capability. Communication noise increases when transfer efficiency is low. The received peer information is therefore both relationship-dependent and imperfect.

\subsection{Hybrid support}
Hybrid support integrates AI and peer information. Integration ability is
\begin{equation}
I_i^{H}=\clip(0.18+0.28v_i+0.26a_i+0.16s_i+0.12V,0.10,0.96).
\end{equation}
The integrated source vector interpolates between the average of AI and peer information and the dimension-wise better source. Disagreement is penalized by
\begin{equation}
0.20(1-I_i^{H})\lvert\mathbf g_i^{AI}-\mathbf g_{ij}^{P}\rvert.
\end{equation}
Hybrid coordination burden declines when the task requirement vector is more multidimensional, using normalized entropy of $\mathbf r$.

\subsection{Task quality and support gain}
For temporary effective skill vector $\mathbf{x}$, requirement vector $\mathbf r$, difficulty $d$, confidence $c_i$, and extra burden $b$, latent quality is
\begin{equation}
Q^{*}=\sigm\left(3.35(\mathbf{x}\cdot\mathbf r-0.50)-2.20(d+b-0.50)+0.22(c_i-0.50)\right).
\end{equation}
Gaussian noise with SD 0.024 is added and the final quality is clipped to $[0,1]$. Solo quality $Q_i^{solo}$ is computed from the agent's own skill vector. Marginal support gain is
\begin{equation}
G_i=Q_i-Q_i^{solo}.
\end{equation}
This quantity is used in local observational learning and, for peer interaction, social-tie updating.

\subsection{Experience and trust updates}
For the selected route, each task-anchor expectation is updated by
\begin{equation}
V_{iam}'=V_{iam}+\alpha_a(Q_i-V_{iam}),
\end{equation}
where
\begin{equation}
\alpha_a=0.32(0.25+0.75\omega_a)
\end{equation}
and $\omega_a$ is current task similarity to anchor $a$. AI trust follows
\begin{equation}
t_i^{AI'}=\clip\left[t_i^{AI}+0.12(S_i^{AI}-t_i^{AI}),0.01,0.99\right],
\end{equation}
where $S_i^{AI}$ is the task-relevant AI source quality. Partner-specific peer trust uses the analogous update toward experienced peer source quality. Partner expertise expectations are updated with the same task-similarity-weighted route-learning rate.

\subsection{Knowledge and confidence updates}
All tasks generate a small practice component:
\begin{equation}
\Delta\mathbf{k}_i^{practice}=0.008\,Q_i\,\mathbf r\odot(1-\mathbf k_i).
\end{equation}
Novel information is the positive component of the difference between effective supported capability and the agent's pre-episode skill. Additional source learning is
\begin{equation}
\Delta\mathbf{k}_i^{support}=0.095\,A_i^{abs}\,Q_i\,\mathbf n_i\odot\mathbf r,
\end{equation}
where $A_i^{abs}$ depends on the selected route and agent attributes. Confidence updates slowly toward realized quality at rate 0.030.

\subsection{Network updating}
If peer $j$ contributes marginally to the focal outcome, an existing tie updates according to
\begin{equation}
\Delta w_{ij}=\begin{cases}
0.085\,G_{ij}^{P}(1-w_{ij}), & G_{ij}^{P}\ge 0,\\
0.040\,G_{ij}^{P}w_{ij}, & G_{ij}^{P}<0.
\end{cases}
\end{equation}
A discovered peer with $G_{ij}^{P}>0.010$ can form a new tie initialized at 0.34. Every unused active tie is multiplied by $(1-0.002)$ each episode and removed below 0.050.

\subsection{Local observational learning}
For each focal agent, existing neighbors are observed independently with probability 0.45. For each route and task anchor, observed neighbors' marginal support gains are averaged using tie strength and task similarity as weights. The corresponding local route signal then moves toward this observed value at rate 0.13.

\subsection{Output measures}
The model records:
\begin{itemize}
\item Peer-first, AI-first, and Hybrid shares;
\item human-interaction and AI-use rates;
\item mean task quality and success rate;
\item mean knowledge, knowledge Gini, performance Gini, and initial--final knowledge rank correlation;
\item mean AI trust and mean peer trust;
\item network density, mean weighted degree, mean active-tie strength, recent active-tie fraction, edge-retention ratio, new-edge ratio, and strong-tie fraction;
\item mean unique peer partners;
\item mean marginal peer contribution and AI-availability probability.
\end{itemize}
The final summary uses a tail window of at least 20 episodes, depending on the experiment.

\section{Simulation experiments}
\subsection{Reference system and human-only benchmark}
The AI-enabled reference system and the human-only benchmark each use 30 paired replications with common seed construction. The human-only benchmark sets AI availability to zero while preserving agent, task, peer, learning, and network mechanisms. Two-sided 95\% Student-$t$ confidence intervals summarize replication uncertainty.

\subsection{Task--route mechanism analysis}
Twenty replications record task-level events and agent summaries. Outputs include task-by-route quality, marginal support gain, task-specific choice shares, early-versus-late choice adaptation, requirement-dimension outcomes, and experience-to-next-choice reinforcement.

\subsection{Environmental factorial}
The factorial experiment contains $3\times3\times2\times2=36$ conditions with 20 replications per condition:
\begin{center}
\begin{tabular}{ll}
\toprule Factor & Levels \\
\midrule
Mean AI access & 0.35, 0.60, 0.85 \\
AI reliability & 0.55, 0.74, 0.90 \\
Verification support & 0.30, 0.70 \\
Peer-response probability & 0.60, 0.90 \\
\bottomrule
\end{tabular}
\end{center}
The design estimates how opportunity conditions shape interaction composition, task outcomes, learning, trust, and social relations.

\subsection{AI reliability shock experiment}
Three conditions use 30 replications each. Reliability starts at 0.86. The temporary-shock condition drops to 0.48 halfway through the 150-episode horizon and returns to 0.86 after three quarters of the horizon. The persistent-shock condition drops to 0.48 at mid-horizon and remains there. Mean AI access is 0.70 in the shock experiment. Paired late-window contrasts compare post-recovery behavior to the stable condition.

\subsection{Mechanism sensitivity}
Sixteen behavioral and update parameters are varied by $\pm20\%$ around their reference values using paired seeds and 20 replications per setting:
\begin{itemize}
\item expected-quality weight;
\item local social-learning weight;
\item choice temperature;
\item mode-value learning rate;
\item trust learning rate;
\item knowledge learning rate;
\item tie reinforcement rate;
\item tie decay;
\item external peer-search probability;
\item hybrid coordination burden;
\item task-similarity temperature;
\item peer-expertise visibility;
\item peer communication-noise SD;
\item maximum peer-transfer efficiency;
\item AI ambiguity verification cost;
\item AI access saturation.
\end{itemize}

\section{Implementation and reproducibility}
The model is implemented in Python. The repository contains:
\begin{itemize}
\item \texttt{model.py}: entities, task generation, route choice, outcome, learning, trust, and network mechanisms;
\item \texttt{experiments.py}: reference, task-mechanism, factorial, shock, and sensitivity designs;
\item \texttt{run\_analysis.py}: complete analysis pipeline;
\item \texttt{precheck.py}: lightweight diagnostic run;
\item \texttt{figures.py} and \texttt{regenerate\_figures.py}: figure generation from saved CSV results;
\item \texttt{results/}: compact replication-level and summary outputs;
\item \texttt{docs/}: this ODD protocol, full parameter documentation, and reproducibility instructions.
\end{itemize}

Full analysis is executed with
\begin{verbatim}
python run_analysis.py
\end{verbatim}
and figures can be regenerated with
\begin{verbatim}
python regenerate_figures.py
\end{verbatim}
Experiment functions are checkpoint-aware. Replication seed $r$ is generated as $20260821+1009r$, enabling common-random-number paired comparisons. The large event-level file \texttt{task\_mode\_agent\_events.csv} is generated locally because its full reference-study size exceeds the standard GitHub per-file limit; compact summaries used for the reported results are included in the repository.

\section{Model boundaries and limitations}
The model is designed for mechanism exploration. Parameter values are transparent reference settings rather than empirical estimates. Task dimensions represent abstract requirement components rather than named occupational skills. AI reliability is represented as a controllable technological property, whereas real systems may exhibit domain-specific and temporally correlated error. Peer communication is dyadic and abstract, and institutional constraints such as hierarchy, incentives, workload, and formal team structure are not represented. Agents adapt through simplified reinforcement-like rules rather than richer cognitive models. Social ties represent repeated interaction opportunities and relationship strength, not the full multidimensional meaning of real social relationships.

These boundaries limit direct numerical generalization, but they make the feedback mechanisms inspectable. Empirical extensions can calibrate task profiles, AI performance, peer-transfer efficiency, trust learning, and network dynamics using field or laboratory data.

\section{References}
\begin{thebibliography}{9}
\bibitem{grimm2006standard} Grimm, V., Berger, U., Bastiansen, F., Eliassen, S., Ginot, V., Giske, J., Goss-Custard, J., Grand, T., Heinz, S.K., Huse, G., et al.: A standard protocol for describing individual-based and agent-based models. Ecological Modelling \textbf{198}(1--2), 115--126 (2006)
\bibitem{grimm2010odd} Grimm, V., Berger, U., DeAngelis, D.L., Polhill, J.G., Giske, J., Railsback, S.F.: The ODD protocol: A review and first update. Ecological Modelling \textbf{221}(23), 2760--2768 (2010)
\bibitem{grimm2020odd} Grimm, V., Railsback, S.F., Vincenot, C.E., Berger, U., Gallagher, C., DeAngelis, D.L., Edmonds, B., Ge, J., Giske, J., Groeneveld, J., et al.: The ODD protocol for describing agent-based and other simulation models: A second update to improve clarity, replication, and structural realism. Journal of Artificial Societies and Social Simulation \textbf{23}(2), 7 (2020)
\bibitem{muller2013odd} M\"uller, B., Bohn, F., Dre\ss ler, G., Groeneveld, J., Klassert, C., Martin, R., Schlueter, M., Schulze, J., Weise, H., Schwarz, N.: Describing human decisions in agent-based models -- ODD+D, an extension of the ODD protocol. Environmental Modelling \& Software \textbf{48}, 37--48 (2013)
\bibitem{epstein2006generative} Epstein, J.M.: Generative Social Science: Studies in Agent-Based Computational Modeling. Princeton University Press (2006)
\bibitem{mcfadden1974conditional} McFadden, D.: Conditional logit analysis of qualitative choice behavior. In: Zarembka, P. (ed.) Frontiers in Econometrics, pp. 105--142. Academic Press, New York (1974)
\bibitem{simon1955behavioral} Simon, H.A.: A behavioral model of rational choice. Quarterly Journal of Economics \textbf{69}(1), 99--118 (1955)
\bibitem{watts1998collective} Watts, D.J., Strogatz, S.H.: Collective dynamics of small-world networks. Nature \textbf{393}, 440--442 (1998)
\end{thebibliography}
\end{document}
